% Draft for Section 3 (From a 129-Column Genome Record to an Identifiable CAD)
\section{From a 129-Column Genome Record to an Identifiable CAD}
\label{sec:genome}

\subsection{A Gene Is a Contract, Not Merely a Number}
A gene requires an explicit coordinate index, physical or discrete type, unit, box domain, coupling law, and decoder revision. A serialized floating-point token alone cannot distinguish a millimeter normal displacement from an angular bevel setting, a discrete feature-count selector, or a binary structural-layer switch. Those semantics determine whether a decoded B-rep or finite-element mesh is reproducible across software revisions.

Jarambit's native crescent representation encodes 128 continuous numerical coordinates paired with one categorical alloy state, and decodes them into a planform and 2.5D thickness field. In every evaluated population table across Cohort~3 (\text{Run~2}--\text{Run~11}, $1{,}954{,}560$ evaluations across 10 runs of 801 generations) and Cohort~1 (\text{Run~1}, $1{,}984$ historical births), each row contains 146 columns: 17 scalar evaluation and search-metadata columns (columns $0..16$), followed by 128 numerical \code{f64} gene coordinates (columns $17..144$) and one categorical alloy column (\code{material}, column $145$, taking values \code{Ti-6Al-4V} or \code{Ti-3Al-2.5V}), yielding 129 genome-and-material fields per birth record.

\subsection{The Four Structural Blocks of the 128 Numeric Genes}
The 128 numeric coordinates partition into four functional blocks that reflect the structural hierarchy of the decoder:
\begin{enumerate}[leftmargin=*,nosep]
\item \textbf{Block A: Silhouette Normal Control Points (\code{cp00}--\code{cp39}, indices $0..39$, 40 coordinates).} Each coordinate specifies a normal displacement in millimeters from the 40-station reference outline. Station-specific bounds restrict displacements to $\pm 1.5\,\mathrm{mm}$ on the outer ring arc ($i\in[0,4]\cup[36,39]$), $\pm 2.0\,\mathrm{mm}$ along the tip approach ($i\in[15,22]$), and $\pm 3.5\,\mathrm{mm}$ along the belly, spine, and handle arches.
\item \textbf{Block B: Planform, Stock, Bevel, Relief, Lattice, Fuller, and Edge Coordinates (indices $40..89$, 50 coordinates).} Indices $40..44$ parameterize global planform and ring dimensions (\code{blade\_drop\_deg}, \code{grip\_dx}, \code{ring\_wall}, \code{ring\_id}, \code{neck\_fillet}); indices $45..49$ define stock thickness \code{t\_stock} ($[4.5,6.0]\,\mathrm{mm}$), tip thickness \code{t\_tip} ($[2.2,5.0]\,\mathrm{mm}$), taper origin \code{taper\_x0}, and inside/spine bevel angles (\code{edge\_angle\_in}, \code{edge\_angle\_sp}); indices $50..89$ parameterize the comet groove, ergonomics, three-tier serration depths, handle relief pocket (\code{relief\_d}, \code{relief\_w}, \code{relief\_q}), dimple lattice, blade fuller, handle flute, surface texture mode, cord wrap, grip pad, and cutting-edge thickness \code{edge\_t} ($[0.25,1.20]\,\mathrm{mm}$).
\item \textbf{Block C: Internal-Machining Coordinates (indices $90..116$, 27 coordinates).} These 27 coordinates control triangular isogrid pocketing (\code{iso\_d} through \code{iso\_grade}, \code{iso\_handle\_start}, \code{iso\_land}), pierced skeletonization windows (\code{win\_n} through \code{win\_decay}), ring lightening flutes (\code{rflute\_d}, \code{rflute\_w}, \code{rflute\_arc}), and angled micro-perforation slots (\code{perf\_d} through \code{perf\_handle\_start}).
\item \textbf{Block D: Continuous-Rail Layer (indices $117..127$, 11 coordinates).} The final 11 numeric coordinates encode the binary layer switch \code{p1\_on} ($\{0,1\}$), continuous perimeter rail width \code{rail\_w} ($[2.2,5.5]\,\mathrm{mm}$), solid neck length \code{neck\_l} ($[6.0,24.0]\,\mathrm{mm}$), interior pocket depth and corner radius (\code{pocket\_d}, \code{pocket\_r}), active spline-knot count \code{k\_n} ($\{3,4,5\}$), and five interior thickness knots \code{k0}--\code{k4} ($[1.5,6.0]\,\mathrm{mm}$).
\end{enumerate}

\subsection{Decoder Coupling, Clamping, and Inactive Coordinates}
These 129 fields are not independent degrees of freedom. Before phenotype construction, the decoder clamp enforces box bounds on all 128 numeric coordinates together with deterministic cross-gene coupling laws:
\begin{equation}
 t_{\text{tip}} \le t_{\text{stock}},\quad x_{1,\text{comet}} \ge x_{0,\text{comet}} + 15\,\mathrm{mm},\quad s_{1,\text{fuller}} \ge s_{0,\text{fuller}} + 0.12,\quad s_{1,\text{win}} \ge s_{0,\text{win}} + 0.12,
\end{equation}
while rounding discrete coordinates (\code{texture\_mode} $\in\{0,1,2\}$, \code{win\_n} $\in\{0,\dots,8\}$, \code{p1\_on} $\in\{0,1\}$, and \code{k\_n} $\in\{3,4,5\}$). When \code{p1\_on} $< 0.5$, the continuous-rail thickness law is inactive in the decoder even though its 11 numeric coordinates remain serialized in the genome record; conversely, when \code{p1\_on} $\ge 0.5$, the continuous perimeter rail and spline-knot thickness law supersede legacy interior pocket features. Preserving all 129 columns losslessly is therefore essential to distinguish an inactive coordinate mutation from a geometric change.

\subsection{Candidate, Genome, and Export Identity}
Because multiple birth events can produce identical clamped coordinates (for example, Cohort~1 \text{Run~1} records $1{,}984$ births across $1{,}965$ distinct gene serializations), Jarambit separates birth/individual identity (\text{Individual~ID}), genome serialization identity, population-slot appearance, and exported CAD/assessment bundle identity. For Individual~2000, the STEP solid ($1$ solid, $1$ shell, $55{,}698$ face tokens), the render mesh ($44{,}994$ vertices, $90{,}000$ triangles), the $0.25\,\mathrm{mm}^2$ plate mesh ($13{,}503$ nodes, $25{,}209$ triangles), and the $P_2$ tetrahedral solid tip submodel ($91{,}008$ tetrahedra, $416{,}583$ displacement degrees of freedom) represent distinct discretizations with separate provenance bindings documented in the Supplement.
